% Pursuing Stacks (cotes 134-2 to 134-8): archive cote 134-3, site page 8 (PDF page 9), typescript page 8.
% typescript. Draft compared: ps2.tex:335-384.
% AI-transcribed by Claude Opus 5.5 (claude-opus-5-5, adaptive thinking, effort high), Message Batches API, batch ps1 (msgbatch_01Rv5hCaNTtUt4JHx1yqoHPL), 2026-11-12; not yet checked by eye.
\scanpage{134-3/8}
\margin{8}
\mspage{8}

just localizing (\over{C}{C}at) with respect to weak isomorphisms, namely func-%
tors $C$ \hand{$\to$} $C'$ inducing an isomorphism between the corresponding homoto-%
py types. Pushing a little further in this direction, one may conjectu%
re that (Cat) is a \hand{``}closed model category,\hand{''} whose weak equivalences are
the functors just specified, and whose cofibrations are just functors
which are injective on objects and injective on arrows.\hand{(19).}

The fact that categories are suitable objects for defining \inserted{arbitrary} homotopy
types is quite remarkable, and one moreover it seems which has not found
it's way still into the minds of topologists o\over{f}{r} homotopists, with only
few exceptions I guess. These objects are extremely simple and familiar
to most mathematicians; what is somewhat more sophisticated is the
process of localization towards homotopy types, or equivalently, the
explicit description of weak equivalences, within the framework of
\inserted{usual} category theory. This would amount more or less to the same as descri-%
bing the homotopy groups of a category, which does not seem any simpler
than the same task for its nerve. As for the cohomology invariants, which
can be interpreted as the left derived functors of the $\underset{\hand{$\leftarrow$}}{\lim}_{C}$ functor ,
or rather its values on particular presheaves (for instance constant
presheaves), they are of course known to be significant\hand{,} independently
of any particular topological interpretation, but they are not expressi%
-ble either in direct terms. (The most common computation for these is
again via the nerve of C.)

This situation suggests that for any \xout{w}natural integer n\hand{$\geq$}1, the
category of n-stacks can be used as a category of models for \inserted{the,} usual homo-%
topy category, in particular any n-stack gives rise to a homotopy type,
\xout{whilexthe} and up to equivalence we should get any homotopy type in this
way (for instance, through the n-category canonically associated to
any 1-category giving rise to this homotopy type). The homotopy types
coming from n-Gr-stacks, however, should be merely the n-truncated
ones, namely those whose homotopy groups in dimension \hand{$>$} n are zero.
\xout{Theyxappearx} Moreover, \xout{as} n-Gr-stacks appear as \inserted{([in some respects])} the most ad\unsure{\over{e}{e}}quate \margin{\illegible{1}}
algebraic structures for expressing n-truncated homotopy types, the
latter being deduced from the former, presumably, by the same process
of localization by weak equivalences. Moreover, in the context\deleted{s} of
n-Gr-stacks, the notions of homotopy groups and of weak equivalences
are described in a particularly obvious way. Thus, passing to the
limit case n=$\infty$ , it is $\infty$-Gr-stacks rather than general $\infty$-stacks
which appear as the neatest algebraic description of homotopy types.
